A principal contribuição deste trabalho é uma comparação quantitativa entre o processamento
em borda e em nuvem para a leitura automática de medidores por \gls{ocr}, feita no mesmo
hardware, com os mesmos modelos e no mesmo enlace \gls{lorawan}, de modo que as diferenças
medidas decorrem apenas do local do processamento. Essa comparação atende ao foco declarado
na Seção~\ref{sec:justificativa}, o de oferecer insumos para a tomada de decisões
arquiteturais em sistemas \gls{iot} com redes de baixa largura de banda, e reúne em um mesmo
cenário a acurácia, o custo computacional no nó, a latência, a entrega e o consumo por
leitura.

Para a Ciência da Computação, os resultados trazem quatro contribuições. A primeira diz
respeito à decisão de onde executar a computação. Em redes de longo alcance e baixa taxa, o
custo de comunicação dominou o de computação: mesmo cerca de 1.000 vezes mais lento que o
servidor na leitura dos dígitos, o nó entregou a leitura com latência cerca de 42 vezes
menor e consumo 16,7 vezes menor, porque transmite cerca de 350 vezes menos dados. A
vantagem na autonomia, porém, depende da corrente do nó em repouso, que passa a dominar o
consumo em intervalos longos entre leituras.

A segunda diz respeito ao aprendizado de máquina embarcado. A quantização completa em INT8
após o treinamento não reduziu a acurácia de forma estatisticamente significativa em um
classificador de 81,5~KB (Seção~\ref{subsec:resultado_h1}), sem a necessidade do
treinamento consciente da quantização. O que limitou a acurácia foi a qualidade da entrada,
e não o modelo: com a câmera fixa, uma região de interesse calibrada na instalação e o
envio ao classificador apenas da face de cada roda levaram a 83\% das leituras e a 97\% dos
dígitos corretos no hidrômetro físico.

A terceira diz respeito ao método de avaliação. Além de controlar as variáveis entre os
fluxos, a comparação verificou as diferenças com testes estatísticos exatos, adequados a
amostras pequenas (Seção~\ref{subsubsec:testes_estatisticos}), e avaliou o classificador em
um conjunto de teste mantido fora de todo o treinamento.

Por fim, o trabalho disponibiliza, em repositório público, os firmwares dos dois fluxos, a
ferramenta de calibração da região de interesse, os \textit{scripts} de recepção,
remontagem e medição e os dados das medições, o que permite reproduzir e estender a
comparação. Também documenta uma configuração de enlace \gls{lorawan} no plano AU915,
enquadrada no item 10.3 do Ato n\textordmasculine~14.448 \cite{anatel2017}, que permite
enviar uma imagem em vários pacotes por um canal fixo de 500~kHz, sem salto de frequência.
